% ============================================================
%  ch02c  §2.3 Jacquet--Langlands 提升的几何实现
%  原文对应：ICM-arXiv.tex 1011--1069 行
%  编号：定义 2.3.1；猜想 2.3.2；公式 (2.6)(2.7)(2.8)；定理 2.3.3
% ============================================================

\subsection{Jacquet--Langlands 提升的几何实现}\label{geometric JL}

数域上的整体朗兰兹对应要复杂得多。事实上，理论中有分析的部分，目前看来无法纳入范畴化／去范畴化的框架。即便只限于理论的代数／算术部分，也会有来自 $\ell$ 处与 $\infty$ 处的复杂性。特别是，$\ell$ 处与 $\infty$ 处的局部朗兰兹对应的范畴形式尚未被完全理解。

尽管如此，在与 Emerton 以及 Emerton--Gee 的合作文章（\cite{Emerton-Zhu, Emerton-Gee-Zhu}）中，我们将对志村簇的上同调、乃至一般局部对称空间的上同调，提出猜想性的 Galois 理论描述，与猜想 \ref{C: gen coh} 平行。本小节只回顾 \cite{Zhu-Conjecture} 中关于借助志村簇上同调几何地实现 Jacquet--Langlands 提升的一个猜想，并讨论 \cite{Xiao-Zhu-Cycle, Hemo-Zhu-Unipotent} 中走向该猜想的若干结果。

先固定一些记号与假设，本小节固定一个素数 $p$。以 $\bA_f=\prod'_q\bQ_q$ 记 $\bQ$ 的有限 adèle 环，$\bA_f^p=\prod'_{q\neq p}\bQ_p$。记 $\overline\eta=\Spec \overline{\bQ}$，其中 $\overline{\bQ}$ 是 $\bQ$ 在 $\bC$ 中的代数闭包。对志村数据 $(G,X)$，以 $\mu$ 记由 $(G,X)$ 按通常方式确定的 $\hat{G}$（关于 $(\hat{B},\hat{T})$）的（微小的）支配权，以 $V_\mu$ 记 $\hat{G}$ 的以 $\mu$ 为最高权的微小表示。以 $E\subset\overline{\bQ}\subset\bC$ 记 $(G,X)$ 的反射域，并记 $d_\mu=\dim X$。

对水平（即开紧子群）$K=K_pK^p\subset G(\bQ_p)G(\bA_f^p)$，以 $\Sh_K(G)$ 记相应水平的志村簇（定义在反射域 $E$ 上），以 $\Sh_K(G)_{\overline\eta}$ 记它沿 $E\to \overline\bQ$ 的基变换。设 $v$ 是 $E$ 上位于 $p$ 之上的位置。所谓特化 $\on{sp}: \overline\eta\to \overline{v}$，指从 $\overline\eta$ 到 $\mO_E$ 在 $v$ 处的严格 Hensel 化的态射。

为避免 Galois 上同调带来的诸多复杂情形（例如扩张纯内形式与内形式之别），也避免几何上的一些复杂情形（例如志村簇与 Shtuka 模空间之间的关系），本小节余下部分设 $G$ 为伴随型；一般 $G$ 的情形见 \cite{Xiao-Zhu-Cycle}，对 $G$ 限制更少的版本见 \cite{Zhu-Conjecture}。

\begin{dfn}\label{D: prime-to-p trivial}
设 $G$ 是 $\bQ$ 上的连通约化群。$G$ 的一个 $p$-外（分别地，有限）平凡化内形式，指 $\bQ$ 上的一个 $G$-挠子 $\beta$，并带有 $\beta$ 在 $\bA_f^p$ 上（分别地，在 $\bA_f$ 上）的平凡化。此时 $G':=\Aut(\beta)$ 是 $G$ 的内形式（因此 $G$ 与 $G'$ 的对偶群被典范地等同），并带有同构 $\theta: G(\bA^p_f)\cong G'(\bA^p_f)$（分别地，$\theta: G(\bA_f)\cong G'(\bA_f)$）。
\end{dfn}

今设 $(G,X)$ 与 $(G',X')$ 是两个志村数据，$G'$ 是 $G$ 的 $p$-外平凡化内形式。借助 $\theta$，可把 $K^p\subset G(\bA_f^p)$ 搬为开紧子群 ${K'}^p\subset G'(\bA_f^p)$。我们把 $p$-外（导出）Hecke 代数 $H_{K^p,\La}$（定义方式与 \eqref{E: derived hecke global} 相同）与 $H_{{K'}^p,\La}$ 等同，并统记作 $H_{K^p,\La}$。设 $K'_p\subset G'(\bQ_p)$ 是开紧子群，记 $K'=K'_p{K'}^p$ 为相应水平。

固定 $G_{\bQ_p}$ 与 $G'_{\bQ_p}$ 的一个拟分裂内形式 $G^*_{\bQ_p}$，并带有钉 $(B^*_{\bQ_p},T^*_{\bQ_p},e^*)$；把 $G_{\bQ_p}$ 实现为 $J_b$、把 $G'_{\bQ_p}$ 实现为 $J_{b'}$，其中 $b,b'\in B(G^*_{\bQ_p})$。在 $G$ 与 $G'$ 为伴随型的假设下，这样的 $b,b'$ 存在且唯一。于是按 \eqref{E: sheaf A}，在 $G^*_{\bQ_p}$ 于 $\La$ 上的局部朗兰兹参数叠 $\Loc_{\CG,p}\otimes\La$ 上有猜想凝聚层 $\frakA_{K_p,\La}$ 与 $\frakA_{K'_p,\La}$。

\begin{conj}\label{C: geo JL}
对特化映射 $\on{sp}: \overline{\eta}\to \overline{v}$ 的任一选取，存在自然映射
\begin{multline}\label{E: geoJL}
R\Hom_{\Coh(\Loc_{\CG,p}\otimes\La)}\bigl(\widetilde{V_\mu}\otimes \frakA_{K_p,\La}, \widetilde{V_{\mu'}}\otimes \frakA_{K'_p,\La}\bigr)\\
\to R\Hom_{H_{K^p,\La}}\bigl(C_c(\Sh_{K}(G)_{\overline\eta}, \La[d_{\mu}]), C_c(\Sh_{K'}(G')_{\overline\eta},\La[d_{\mu'}])\bigr),
\end{multline}
并与复合相容。特别地，有（$E_1$-）代数同态
\begin{equation}\label{E: S operator}
S:R\End_{\Coh(\Loc_{\CG,p}\otimes\La)}\bigl( \widetilde{V_\mu}\otimes \frakA_{K_p,\La}\bigr)\to R\End_{H_{K^p,\La}}\bigl(C_c(\Sh_{K}(G)_{\overline\eta},\La)\bigr),
\end{equation}
它与 \eqref{E: geoJL} 相容。此外，诱导出的作用
\begin{equation}\label{E: S action}
H_{K_p,\La}\stackrel{\eqref{E: Hecke alg isom}}{\cong} R\End(\frakA_{K_p,\La})\to R\End\bigl(\widetilde{V_\mu}\otimes \frakA_{K_p,\La}\bigr)\xrightarrow{S} R\End_{H_{K^p,\La}}\bigl(C_c(\Sh_{K}(G)_{\overline\eta},\La)\bigr)
\end{equation}
与 $H_{K_p,\La}$ 在 $C_c(\Sh_{K}(G)_{\overline\eta},\La)$ 上的自然 Hecke 作用一致（因此与特化映射 $\on{sp}$ 无关）。
\end{conj}

如果对 $C_c(\Sh_{K}(G)_{\overline\eta},\La)$ 有类似于猜想 \ref{C: gen coh} 的 Galois 理论描述，本猜想便是其推论；但本猜想的表述并不需要 $\bQ$ 的整体朗兰兹参数叠存在。无论如何，走向 $C_c(\Sh_{K}(G)_{\overline\eta},\La)$ 的 Galois 理论描述，或许需要以猜想 \ref{C: geo JL} 为输入。还要指出，与函数域情形一样，\cite[Sect. 4.7]{Zhu-Conjecture} 中有本猜想的更一般版本，允许“广义水平结构”，从而使 Igusa 簇的上同调也能出现。

下述定理在特殊情形下验证了该猜想。

\begin{thm}
\label{introT: spectral action}
设志村数据 $(G,X)$ 与 $(G',X')$ 都是阿贝尔型，且 $G'$ 是 $G$ 的有限平凡化内形式。设 $G_{\bQ_p}$ 非分歧（从而 $G'_{\bQ_p}$ 也非分歧）。
\begin{enumerate}
\item\label{introT: spectral action-1} 当 $\La=\bQ_\ell$ 且 $K_p\subset G(\bQ_p)$、$K'_p\subset G'(\bQ_p)$ 为抛孔子群（Bruhat--Tits 意义下）时，映射 \eqref{E: geoJL}（从而 \eqref{E: S operator}）存在。
\item\label{introT: spectral action-2} 若 $K_p$ 超特殊，则当 $\La=\bZ_\ell$ 时映射 \eqref{E: geoJL}（从而 \eqref{E: S operator}）存在，至少对非导出的 $\Hom$ 空间如此。此外，$H_{K_p}^{\cl}$ 通过 \eqref{E: S action} 在 $H_c^*(\Sh_{K}(G)_{\overline\eta},\La)$ 上的作用与 $H_{K_p}^{\cl}$ 的自然作用一致。
\end{enumerate}
\end{thm}

第 \eqref{introT: spectral action-1} 部分在 \cite{Xiao-Zhu-Cycle, Hemo-Zhu-Unipotent} 中证明。证明包含两个要素：一是 \cite{Xiao-Zhu-Cycle} 对 $\Sh_K(G)$ 与 $\Sh_{K'}(G')$ 的 mod $p$ 纤维之间物理对应的构造（这里用到了目前需要假设 $G$ 与 $G'$ 在 $p$ 处非分歧这一点）；二是定理 \ref{T: tame local Langlands}（因此需要 $\La=\bQ_\ell$）。当 $K_p$ 超特殊时，可以取 $\bZ_\ell$ 系数，因为 \eqref{E: local S-operator} 的（非导出版本）对 $\bZ_\ell$ 系数成立，这得益于 \cite{Yu-Jacquet-Langlands-Higher-Weights}。事实上，此时还可以允许志村簇上有非平凡局部系统（见 \cite{Yu-Jacquet-Langlands-Higher-Weights}）。

最后一个断言即志村簇情形的 $S=T$。当 $d_{\mu}=\dim\Sh_K(G)=0$ 时它包含于 \cite{Xiao-Zhu-Cycle}；一般情形在 \cite{Wu-S=T, Zhu-S=T} 中借助 \cite{Scholze-Berkeley, Fargues-Scholze-Geometrization} 的基础性工作得到证明。
